% =============================================================================
% editions.tex — Edition-gated environments for Learning to Lens
% =============================================================================
% \input by main.tex AFTER the \ifsolutions / \ifannotated switches are set and
% AFTER % =============================================================================
% preamble.tex — Shared LaTeX preamble for Learning to Lens notes
% =============================================================================
% This file is \input{} by main.tex and provides all packages, macros, and
% formatting settings used across the tutorial modules.
% =============================================================================

% --- Core math and physics packages ---
\usepackage{amsmath,amssymb,amsthm}
\usepackage{mathtools}          % Extensions to amsmath (e.g., \coloneqq)
\usepackage{physics}            % Convenient notation: \dv, \pdv, \grad, \curl, etc.
\usepackage{tensor}             % Index notation: \indices{^a_b}
\usepackage{bm}                 % Bold math symbols

% --- Figures and graphics ---
\usepackage{graphicx}
\usepackage{float}              % [H] placement for figures
\usepackage{subcaption}         % Subfigures

% --- Tables ---
\usepackage{booktabs}           % Professional tables (\toprule, \midrule, \bottomrule)
\usepackage{array}

% --- Typography and layout ---
\usepackage[utf8]{inputenc}
\usepackage[T1]{fontenc}
\usepackage{microtype}          % Improved typesetting
\usepackage{enumitem}           % Customizable lists
\usepackage{fancyhdr}           % Headers and footers
\usepackage{titlesec}           % Section formatting

% --- Colors and links ---
\usepackage[dvipsnames]{xcolor}
\usepackage[
    colorlinks=true,
    linkcolor=NavyBlue,
    citecolor=ForestGreen,
    urlcolor=BrickRed,
    filecolor=Plum
]{hyperref}

% --- Boxed environments (Part V key results, GR reminders, annotated edition) ---
\usepackage{tcolorbox}          % loads tikz, environ, etoolbox, verbatim
\tcbuselibrary{breakable,skins}
\usepackage{comment}            % \excludecomment: zero-output environments in non-annotated editions (editions.tex)

% --- Code listings (for Mathematica snippets) ---
\usepackage{listings}
\lstdefinestyle{mathematica}{
    language=Mathematica,
    basicstyle=\ttfamily\small,
    keywordstyle=\color{NavyBlue}\bfseries,
    commentstyle=\color{Gray}\itshape,
    stringstyle=\color{BrickRed},
    frame=single,
    breaklines=true,
    numbers=left,
    numberstyle=\tiny\color{Gray},
    backgroundcolor=\color{gray!5},
    captionpos=b
}

% --- Bibliography ---
\usepackage[numbers,sort&compress]{natbib}

% --- Theorem environments ---
\newtheorem{theorem}{Theorem}[section]
\newtheorem{lemma}[theorem]{Lemma}
\newtheorem{proposition}[theorem]{Proposition}
\newtheorem{corollary}[theorem]{Corollary}
\theoremstyle{definition}
\newtheorem{definition}{Definition}[section]
\newtheorem{example}{Example}[section]
\newtheorem{exercise}{Exercise}[section]
\theoremstyle{remark}
\newtheorem{remark}{Remark}[section]
\newtheorem*{note}{Note}

% =============================================================================
% Custom macros — Gravitational lensing notation
% =============================================================================
% Following the conventions of Narayan & Bartelmann (1997),
% Congdon & Keeton (2018), and Saha et al. (2024).
% =============================================================================

% --- Vector quantities (bold, with arrows as needed) ---
\newcommand{\vectheta}{\boldsymbol{\theta}}     % Image position vector
\newcommand{\vecbeta}{\boldsymbol{\beta}}        % Source position vector
\newcommand{\vecalpha}{\boldsymbol{\alpha}}      % Deflection angle vector
\newcommand{\vecxi}{\boldsymbol{\xi}}            % Position in lens plane (physical)
\newcommand{\veceta}{\boldsymbol{\eta}}          % Position in source plane (physical)

% --- Angular diameter distances ---
\newcommand{\Dd}{D_{\mathrm{d}}}                 % Observer-lens distance
\newcommand{\Ds}{D_{\mathrm{s}}}                 % Observer-source distance
\newcommand{\Dds}{D_{\mathrm{ds}}}               % Lens-source distance
\newcommand{\Deff}{D_{\mathrm{eff}}}             % Effective lensing distance

% --- Critical quantities ---
\newcommand{\Sigmacr}{\Sigma_{\mathrm{cr}}}      % Critical surface mass density
\newcommand{\thetaE}{\theta_{\mathrm{E}}}        % Einstein radius (angular)
\newcommand{\RE}{R_{\mathrm{E}}}                 % Einstein radius (physical)
\newcommand{\RS}{R_{\mathrm{S}}}                 % Schwarzschild radius

% --- Lensing quantities ---
\newcommand{\kappab}{\kappa}                     % Convergence
\newcommand{\gammaL}{\gamma}                     % Shear magnitude
\newcommand{\psiL}{\psi}                         % Lensing potential
\newcommand{\alphahat}{\hat{\alpha}}             % Physical deflection angle

% --- Jacobian / magnification ---
\newcommand{\Amat}{\mathcal{A}}                  % Jacobian matrix (inverse magnification)
\newcommand{\Mmat}{\mathcal{M}}                  % Magnification tensor

% --- Cosmology ---
\newcommand{\Omegam}{\Omega_{\mathrm{m}}}        % Matter density parameter
\newcommand{\OmegaL}{\Omega_{\Lambda}}           % Dark energy density parameter
\newcommand{\Hubble}{H_0}                        % Hubble constant

% --- Physical constants ---
\newcommand{\speedoflight}{c}                     % Speed of light
\newcommand{\Gnewton}{G}                         % Newton's gravitational constant
\newcommand{\Msun}{M_{\odot}}                    % Solar mass
\newcommand{\Rsun}{R_{\odot}}                    % Solar radius

% --- Metric and tensor notation ---
\newcommand{\metric}{g_{\mu\nu}}                 % Metric tensor
\newcommand{\christoffel}[3]{\Gamma^{#1}_{\ #2 #3}} % Christoffel symbol
\newcommand{\riemann}[4]{R^{#1}_{\ #2 #3 #4}}  % Riemann tensor
\newcommand{\ricci}{R_{\mu\nu}}                  % Ricci tensor
\newcommand{\ricciS}{R}                          % Ricci scalar

% --- Commonly used shorthands ---
\newcommand{\half}{\tfrac{1}{2}}
\newcommand{\inv}[1]{{#1}^{-1}}
\newcommand{\ee}[1]{\times 10^{#1}}             % Scientific notation: 3\ee{8}
\renewcommand{\order}[1]{\mathcal{O}\!\left(#1\right)} % Order of magnitude (override physics pkg)

% --- Cross-references to Mathematica files ---
% Links point to the GitHub repository so they work from any PDF viewer.
% To open files locally, use: ./open_notebook.sh <name>
% Usage: \mathlink{01a_Special_Relativity/lorentz_transforms.wl}{Lorentz transform script}
\newcommand{\mathlink}[2]{\href{https://github.com/rodelcr/Learning_to_Lens/blob/main/Mathematica/#1}{\texttt{#2}}}
% Usage: \solnlink{01a_Special_Relativity/problems_01a.wl}{Solutions script}
\newcommand{\solnlink}[2]{\href{https://github.com/rodelcr/Learning_to_Lens/blob/main/Solutions/#1}{\texttt{#2}}}

% =============================================================================
% Part V macros — compound lenses, DSPL, central point masses
% =============================================================================
% Ported 2026-10-14 from AGEL/202607_Theoretical_spiral_compound/theory/
% preamble_theory.tex. Everything defined above (\vectheta, \Dd, \thetaE,
% \kappab, \Amat, \Omegam, ...) is shared with that preamble by construction and
% is NOT redefined here. \dd is NOT defined here either: the physics package
% (loaded above) already provides \dd = \mathrm{d}; note that physics parses
% \dd(x), \dd[n]{x} as delimited forms — write \dd\theta, \dd^2\xi', never \dd(...).
% =============================================================================
\newcommand{\Brat}{\mathcal{B}}          % DSPL distance ratio  D_{ls1} D_{s2} / (D_{s1} D_{ls2})  (< 1)
\newcommand{\gp}{\gamma'}                % total (3-D) density slope, rho ~ r^{-gamma'}
\newcommand{\thetaBH}{\theta_{\bullet}}  % point-mass (SMBH) Einstein radius
\newcommand{\mBH}{m}                     % central-mass fraction  M_BH / M(<theta_E)
\newcommand{\etamea}{\eta}               % measured Einstein-radius ratio  theta_E1 / theta_E2
\newcommand{\ysrc}{y}                    % signed 1-D source position
\newcommand{\Dls}[2]{D_{#1#2}}           % D_{ls1}-style angular-diameter distance between two planes
\newcommand{\deriv}[2]{\frac{\mathrm{d} #1}{\mathrm{d} #2}}
\newcommand{\pderiv}[2]{\frac{\partial #1}{\partial #2}}

% --- Boxed content environments (ALWAYS visible, all editions) ---
% \begin{keybox}[Title] ... \end{keybox}     key result
% \begin{aside}[Title]  ... \end{aside}      "GR reminder" for Carroll-level readers
\newtcolorbox{keybox}[1][]{enhanced,breakable,colback=Yellow!8,colframe=BrickRed!70,
    title={#1},fonttitle=\bfseries}
\newtcolorbox{aside}[1][]{enhanced,breakable,colback=NavyBlue!4,colframe=NavyBlue!50,
    fonttitle=\bfseries,title={GR reminder\ifblank{#1}{}{ --- #1}}}
 (needs tcolorbox, etoolbox, comment).
%
% Three editions are built by build.sh:
%   student     : solutions OFF, annotations OFF
%   instructor  : solutions ON,  annotations OFF
%   annotated   : solutions ON,  annotations ON   ("Annotated Study Edition")
%
% In the student and instructor editions every environment below is a
% comment-package exclusion: the body is skipped at the CHARACTER level and
% produces no tokens at all, so those two PDFs are content-identical with or
% without annotations in the sources (build.sh asserts this via text-md5
% fingerprints and the uppercase title sentinels ANNOTATION/DERIVATION/BRIDGE).
%
% RULES FOR AUTHORS (linted; see VERIFICATION.md "Annotated Study Edition")
%  R1  \begin{annotation}[title] and \end{annotation} each on their OWN line;
%      \end{...} with NO leading whitespace; nothing else on either line.
%      (comment.sty scans line by line: text after \begin{..} on the same line is
%      swallowed, and an indented \end{..} is not recognised.)
%  R2  Never inside align / gather / multline / equation / tabular / itemize.
%      Boxes sit between paragraphs only. Inside a display use R4.
%  R3  No \label inside an annotated block unless it is only \ref'd from other
%      annotated blocks; never \cite a key that exists only for annotations.
%  R4  Annotated-only material inside a DISPLAY (align/gather/equation):
%      a `\\` inside a skipped \ifannotated branch within an alignment breaks
%      the build ("Incomplete \iffalse"), so extra rows CANNOT be gated row-wise.
%      Two supported mechanisms (both verified by the Phase-0 smoke test):
%      (a) annotated-only prose between rows:  \ifannotated\intertext{...}\fi
%          placed IMMEDIATELY AFTER the row's \\ (amsmath rule; placing it before
%          the \\ yields an extra empty numbered row in the annotated edition)
%      (b) extra rows: duplicate the whole display,
%            \ifannotated
%              \begin{align} ... \notag \\ <extra row> \notag \\ ... \label{eq:x} \end{align}
%            \else
%              \begin{align} ... \label{eq:x} \end{align}
%            \fi
%          with IDENTICAL labels and identical numbered rows in both branches
%          (extra rows \notag) so equation numbers agree across editions.
%  R5  Every numeric claim made inside an annotation gets a check[] in the
%      module's Solutions/NN/problems_NN.wl (Mathematica is the source of truth).
% =============================================================================

\ifannotated
  \ifsolutions\else
    \errmessage{editions.tex: the annotated edition requires solutions (set \showsolutions)}%
  \fi
  \tcbset{ltlbox/.style={enhanced,breakable,left=5pt,right=5pt,top=4pt,bottom=4pt,
          before skip=8pt plus 2pt,after skip=8pt plus 2pt,fonttitle=\bfseries\small}}
  % Commentary on the surrounding text: intuition, unit checks, sign conventions,
  % common errors, "why this matters".
  \newtcolorbox{annotation}[1][]{ltlbox,colback=Plum!4,colframe=Plum!55,
          title={ANNOTATION\ifblank{#1}{}{ --- #1}}}
  % Step-by-step expansion of a derivation the main text compresses.
  \newtcolorbox{derivation}[1][]{ltlbox,colback=gray!6,colframe=gray!60,
          title={DERIVATION\ifblank{#1}{}{ --- #1}}}
  % Forward link: where this result is used in Part V (compound lenses).
  \newtcolorbox{bridge}[1][]{ltlbox,colback=ForestGreen!5,colframe=ForestGreen!55,
          title={BRIDGE\ifblank{#1}{}{ --- #1}}}
  % Short inline aside (TEXT MODE ONLY, never inside math).
  \newcommand{\annot}[1]{{\color{Plum}\small\itshape[#1]}}
\else
  \excludecomment{annotation}
  \excludecomment{derivation}
  \excludecomment{bridge}
  \newcommand{\annot}[1]{\unskip}
\fi
